%!TEX root = ../main.tex
% Appendix

%% ============================================================
\section{Evaluation Dimensions: Detailed Definitions}
\label{app:eval_dimensions}

The \medskillbench evaluation framework scores each skill--task trial along eight binary dimensions.
The first three are designated \emph{core} dimensions; a skill passes overall only if all three core dimensions score~1.
The remaining five are \emph{auxiliary} dimensions that capture execution quality beyond the minimum pass threshold.
Below we provide the full definition and scoring criteria for each dimension.

\paragraph{1.\ \texttt{skill\_invoked} (Core).}
The agent explicitly references, reads, or attempts to use the \skillmd file during the execution trajectory.
A score of~1 is assigned if the agent's trajectory contains evidence that the skill specification was consulted---for example, a file-read action targeting the \skillmd, direct quotation of skill instructions, or an explicit statement acknowledging the skill.
A score of~0 indicates the agent completed the task without any evidence of engaging with the provided skill.

\paragraph{2.\ \texttt{skill\_used\_correctly} (Core).}
The agent applies the skill in a manner consistent with the specification's intended use case, input--output contract, and procedural instructions.
A score of~1 requires that the agent's actions align with the skill's documented workflow---correct parameter usage, appropriate invocation context, and adherence to any constraints or prerequisites stated in the specification.
Partial or superficial references to the skill that do not translate into specification-aligned actions receive a score of~0.

\paragraph{3.\ \texttt{task\_completion} (Core).}
The agent produces a final output that satisfies the task objective.
Scoring considers whether the output addresses the stated goal, provides a substantive answer or artifact, and is not truncated, empty, or obviously incorrect.
A score of~1 does not require perfect accuracy but demands that the output constitutes a reasonable, complete attempt.

\paragraph{4.\ \texttt{tool\_execution\_success} (Auxiliary).}
At least one tool invocation during the trajectory produces a verified, non-trivial output.
This dimension captures whether the agent successfully grounds the skill specification into executable tool calls that return meaningful results.
A score of~0 indicates that all tool calls either failed, produced empty outputs, or were never attempted despite the skill specifying tool-based workflows.

\paragraph{5.\ \texttt{output\_quality} (Auxiliary).}
The final output demonstrates domain-appropriate quality: accurate terminology, logically coherent reasoning, and sufficient depth for the task's complexity.
The judge assesses whether the output would be considered adequate by a domain practitioner, accounting for the difficulty of the generated task.

\paragraph{6.\ \texttt{reasoning\_quality} (Auxiliary).}
The agent's reasoning trajectory---including intermediate steps, tool-call rationale, and error recovery---demonstrates logical coherence and appropriate problem decomposition.
A score of~1 indicates that the agent's chain of thought is traceable, non-circular, and reflects genuine engagement with the problem rather than superficial pattern matching.

\paragraph{7.\ \texttt{efficiency} (Auxiliary).}
The agent completes the task within a reasonable number of iterations relative to task complexity, without excessive redundant actions, repeated failures on the same step, or unnecessary detours.
A score of~0 is assigned when the agent exhausts the iteration budget, enters retry loops without progress, or performs substantially more actions than necessary.

\paragraph{8.\ \texttt{trajectory\_quality} (Auxiliary).}
The overall trajectory exhibits purposeful progression toward the goal, appropriate error handling, and coherent state management across iterations.
This holistic dimension captures trajectory-level phenomena not addressed by individual dimensions: does the agent recover from failures constructively, maintain context across steps, and converge toward a solution?

\paragraph{Aggregation.}
A skill achieves \emph{pass} status when all three core dimensions score~1.
A skill achieves \emph{gold} status when all eight dimensions score~1.
\emph{Silver} status denotes core pass with at least one auxiliary failure.
The \rbu gap is computed as the difference between \texttt{skill\_invoked} and \texttt{task\_completion} (see Equation~\ref{eq:rbu} in the main text).


%% ============================================================
\section{Mandatory Read-First Prompt}
\label{app:read_first_prompt}

Every skill evaluation trial injects the following system-level prompt before the task description to ensure the agent engages with the \skillmd file.
This prompt is appended to the agent's system message and is not visible to the judge.

\begin{small}
\begin{verbatim}
IMPORTANT: Before attempting this task, you MUST
read the file SKILL.md in your current working
directory. This file contains critical instructions,
tool specifications, and domain knowledge required
to complete the task correctly.

Steps:
1. Read SKILL.md completely before taking any action.
2. Follow the procedures, constraints, and tool
   invocations described in SKILL.md.
3. If SKILL.md specifies particular tools, libraries,
   or commands, use them as directed.
4. Only after reading and understanding SKILL.md
   should you proceed to address the task.

Do NOT skip reading SKILL.md. Do NOT rely solely on
your internal knowledge if SKILL.md provides
specific procedures.
\end{verbatim}
\end{small}

This prompt achieves a 96.2\% invocation rate across the full corpus of 1,462 skills, validating its effectiveness at directing agent attention.
The remaining 3.8\% of non-invocation cases arise primarily from skills whose \skillmd files are empty, malformed, or contain only metadata headers without actionable content.


%% ============================================================
\section{Skill Domain Distribution}
\label{app:domain_distribution}

The 917 Tier-A skills in \medskillbench span multiple medical and biomedical domains, reflecting the breadth of open-source medical agent tool repositories.
Domain annotations are produced by an automated classifier with manual verification of borderline cases.

\paragraph{Primary domains.}
The corpus covers the following major areas:
\begin{itemize}[nosep,leftmargin=*]
  \item \textbf{Biology \& Bioinformatics} --- Sequence analysis (BLAST, multiple sequence alignment), gene expression analysis, phylogenetics, genome annotation, protein structure prediction, and pathway analysis tools.
  \item \textbf{Clinical Medicine} --- Clinical decision support, diagnostic reasoning guides, treatment protocol references, drug interaction checkers, and clinical calculator wrappers.
  \item \textbf{Drug Development \& Pharmacology} --- Molecular docking, ADMET prediction, compound screening, pharmacokinetic modeling, and drug--target interaction databases.
  \item \textbf{Chemoinformatics} --- Chemical structure manipulation (RDKit-based), molecular fingerprinting, SMILES/InChI conversion, and chemical property prediction.
  \item \textbf{Proteomics \& Structural Biology} --- Protein--protein interaction analysis, mass spectrometry data processing, structural alignment, and binding site prediction.
  \item \textbf{Genomics \& Variant Analysis} --- Variant calling, VCF processing, genome-wide association analysis, and annotation pipelines.
  \item \textbf{Medical Literature \& Evidence} --- PubMed search, systematic review automation, evidence grading, and citation management.
  \item \textbf{Laboratory \& Safety} --- Laboratory protocol guides, biosafety references, equipment operation procedures, and regulatory compliance checklists.
  \item \textbf{Medical Imaging} --- DICOM processing, radiological analysis aids, and image segmentation tool wrappers.
  \item \textbf{Public Health \& Epidemiology} --- Epidemiological modeling, outbreak analysis, surveillance data processing, and population health statistics.
\end{itemize}

\paragraph{Distribution characteristics.}
The domain distribution is intentionally unbalanced, reflecting the natural distribution of skills in the source repositories.
Biology \& Bioinformatics and Clinical Medicine together account for the largest share, followed by Drug Development and Medical Literature tools.
This imbalance is preserved rather than artificially balanced, as it reflects the real-world landscape of available medical agent tools.


%% ============================================================
\section{Full Benchmark Statistics}
\label{app:full_benchmark}

\Cref{tab:full_benchmark} extends the summary in Table~\ref{tab:benchmark_match} with additional columns for absolute match counts, unique skill counts, and the top match types observed per benchmark.

\begin{table*}[t]
\centering\small
\caption{Extended skill--benchmark matching statistics across 12 medical QA benchmarks (917 Tier-A skills).
``Matched'' indicates the number of questions with at least one skill match.
``Top match types'' lists the two most frequent match type annotations for each benchmark.}
\begin{tabular}{lrrrrll}
\toprule
\textbf{Benchmark} & \textbf{Total Q} & \textbf{Matched} & \textbf{Match\%} & \textbf{Skills} & \textbf{Top Match Type 1} & \textbf{Top Match Type 2} \\
\midrule
Lab-Bench         &   900 & 416 & 46.2 & 90 & \texttt{database\_query}    & \texttt{literature\_search} \\
MedAgents         &   332 &  96 & 28.9 & 16 & \texttt{literature\_search}  & \texttt{reasoning\_guide}   \\
MedMCQA           & 4,183 &1,071& 25.6 & 50 & \texttt{literature\_search}  & \texttt{database\_query}    \\
PubMedQA          & 1,000 & 194 & 19.4 & 35 & \texttt{literature\_search}  & \texttt{database\_query}    \\
HeadQA            & 2,742 & 499 & 18.2 & 66 & \texttt{literature\_search}  & \texttt{reasoning\_guide}   \\
SGI-Reasoning     &    72 &   9 & 12.5 & 17 & \texttt{reasoning\_guide}    & \texttt{literature\_search} \\
MMLU-Med          & 1,089 & 127 & 11.7 & 42 & \texttt{literature\_search}  & \texttt{reasoning\_guide}   \\
MedQA-USMLE       & 1,273 &  71 &  5.6 & 21 & \texttt{literature\_search}  & \texttt{reasoning\_guide}   \\
LabSafety         &   765 &  42 &  5.5 &  3 & \texttt{literature\_search}  & \texttt{reasoning\_guide}   \\
MedCalc-Bench     & 1,067 &  47 &  4.4 &  5 & \texttt{database\_query}     & \texttt{reasoning\_guide}   \\
MedBullets-5op    &   298 &  15 &  5.0 &  6 & \texttt{literature\_search}  & \texttt{reasoning\_guide}   \\
MedBullets-4op    &   298 &   8 &  2.7 &  6 & \texttt{literature\_search}  & \texttt{reasoning\_guide}   \\
\midrule
\textbf{Overall}  & \textbf{14,019} & \textbf{2,596} & \textbf{18.5} & \textbf{178} & \texttt{literature\_search} (71.7\%) & \texttt{database\_query} (21.0\%) \\
\bottomrule
\end{tabular}
\label{tab:full_benchmark}
\end{table*}

\paragraph{Notes on matched counts.}
The ``Matched'' column reports unique questions with at least one skill assignment; a single question may be matched to multiple skills, yielding approximately 3,898 total skill--question pairs across the 2,596 matched questions.
Relevance annotations indicate that 91.5\% of matches are classified as \emph{weak} relevance (the skill provides tangential support) and 8.5\% as \emph{strong} relevance (the skill directly addresses the core reasoning required by the question).

\paragraph{Skill concentration analysis.}
Across all benchmarks, the top three matched skills are \texttt{agent-browser} (appearing in 78\% of all matched pairs), \texttt{ai-analyzer} (13.8\%), and \texttt{bio-entrez-search} (12.2\%).
This heavy concentration underscores that current medical skill repositories favor general-purpose retrieval and analysis tools over domain-specialized reasoning aids.
Lab-Bench is an exception, with 90 unique skills matched---the highest skill diversity of any benchmark---consistent with its emphasis on specific laboratory protocols and bioinformatics pipelines.


%% ============================================================
\section{Judge Calibration: v2 vs.\ v3 Comparison}
\label{app:judge_calibration}

During development, we identified systematic false negatives in our v2 judge that led to an artificially low pass rate of 18.5\%.
The v3 judge, deployed for all results reported in the main text, addresses three calibration issues:

\begin{enumerate}[nosep,leftmargin=*]
  \item \textbf{Partial completion credit}: The v2 judge applied an overly strict interpretation of \texttt{task\_completion}, requiring exact output format matches. The v3 judge evaluates substantive completion regardless of formatting variations.
  \item \textbf{Implicit skill usage}: The v2 judge required explicit mention of the \skillmd file name for \texttt{skill\_used\_correctly}. The v3 judge accepts behavioral evidence of skill adherence even without verbatim references.
  \item \textbf{Tool failure tolerance}: The v2 judge penalized \texttt{task\_completion} when intermediate tool calls failed, even if the agent recovered and produced a correct final output. The v3 judge evaluates the final output independently of intermediate failures.
\end{enumerate}

The resulting pass rate shift from 18.5\% (v2) to 64.4\% (v3) reflects these scoring calibration improvements and does not indicate changes to the underlying skill corpus or execution environment.
All comparative experiments in this paper use the v3 judge exclusively.


%% ============================================================
\section{\gepamed Detailed Effectiveness}
\label{app:gepa_results}

\begin{table}[t]
\centering\small
\caption{\gepamed vs.\ four description-refinement baselines on the high-value skill subset.
Metrics: Core Pass Rate, \rbu Gap, Output Quality (auxiliary dim).}
\begin{tabular}{lccc}
\toprule
\textbf{Method} & \textbf{Pass\%} & \textbf{\rbu Gap} & \textbf{OutQ} \\
\midrule
Original Skill           & 62.3 & 0.215 & 0.824 \\
EASYTOOL-style           & 67.1 & 0.182 & 0.831 \\
PLAY2PROMPT-style        & 69.4 & 0.170 & 0.836 \\
Trace-Free+              & 72.0 & 0.149 & 0.842 \\
\midrule
\textbf{\gepamed (ours)} & \textbf{78.4} & \textbf{0.112} & \textbf{0.857} \\
\bottomrule
\end{tabular}
\label{tab:gepa_results}
\end{table}

\todo{Replace with measured numbers once \gepamed experiments complete.}


%% ============================================================
\section{Batch-Level Evaluation Stability}
\label{app:batch_stability}

To assess evaluation stability across the corpus, we report results stratified by processing batch in \Cref{tab:batch_stability}.
Skills were processed in sequential batches of 400 (with the final batch containing the remaining 262 skills).

\begin{table}[t]
\centering\small
\caption{Evaluation results by processing batch, demonstrating stability across the corpus.}
\begin{tabular}{lrccr}
\toprule
\textbf{Batch} & \textbf{$n$} & \textbf{Pass\%} & \textbf{Gold\%} & \textbf{\rbu Gap} \\
\midrule
1 (1--400)       & 400 & 57.8 & 15.0 & 0.333 \\
2 (401--800)     & 400 & 65.5 & 23.8 & 0.173 \\
3 (801--1200)    & 400 & 68.3 & 22.8 & 0.168 \\
4 (1201--1462)   & 262 & 67.2 & 20.2 & 0.199 \\
\midrule
\textbf{Overall} & \textbf{1,462} & \textbf{64.4} & \textbf{20.5} & \textbf{0.220} \\
\bottomrule
\end{tabular}
\label{tab:batch_stability}
\end{table}

Batch~1 exhibits a lower pass rate (57.8\%) and higher \rbu gap (0.333) than subsequent batches.
This difference is attributable to minor judge prompt refinements applied after Batch~1 processing (within the v3 protocol), which improved consistency in borderline cases.
Batches 2--4 show stable performance, with pass rates in the 65.5--68.3\% range and \rbu gaps between 0.168 and 0.199, confirming that the evaluation framework yields consistent results once calibrated.
